% GENERATED FILE. Edit canonical lessons, then run npm run generate:books.
% canonical-source-hash: fnv1a64:50fb1f567fe4106e
% canonical-lessons: GU-C36-te, GU-C36-ame-aapne, GU-C36-te-write, GU-C36-je, GU-C36-je-te

\chapter{He, She, We, and the One Who}
\label{ch:gu-pronouns-and-relative}
\hlchaptermodality{40}

\begin{chapteropening}
\textbf{By the end of this chapter:} \emph{I can talk about somebody who is not present, choose between the two words for we, and describe a person instead of naming them.}

The absent third person, the two we, and the relative. te covers he, she and it with one word; ame and aapne are the exclusive and inclusive we, where choosing wrongly audibly shuts the listener out; and je ... te ... turns out to be the same two-ended shape as jo ... to ... from the chapter before -- one habit, three uses.
\end{chapteropening}

\section[te]{\gu{તે} (te) — he, she, it}

\label{lesson:GU-C36-te}



\begin{quote}
From the chapter behind you, a condition. And from the third chapter of this book, the word for \textbf{I}.

\begin{itemize}
\raggedright
  \item \emph{Recall:} \emph{jo … to …}, and \emph{hũ}
\end{itemize}

You have been able to talk about yourself since chapter 3 and about the person in front of you since chapter 2. You have never been able to talk about anybody \textbf{absent}.
\end{quote}

\begin{cousinweb}
\begin{quote}
\textbf{\gu{તે}} — \emph{te} — \textbf{he}, \textbf{she}, or \textbf{it}
\end{quote}

One word, three English ones. Gujarati's third person does not make you choose a gender before you speak — which is a real convenience, and occasionally a real ambiguity.

\textbf{\gu{તે}} also means \emph{that}, pointing at something away from both speakers. The pronoun and the pointing word are the same word, which is where the \textbf{\gu{તો}} of the last chapter came from as well.
\end{cousinweb}

\subsection*{Your turn}
\begin{itemize}
\raggedright
  \item \emph{Say it:} \emph{hũ}, then \emph{te}. Yourself, then someone absent.
  \item \emph{Build:} a sentence about someone who is not in the room.
  \item \emph{Say it:} \emph{te} as \emph{that}, pointing at something across the room.
\end{itemize}

\begin{tcolorbox}[breakable,colback=teal!4,colframe=teal!35!black,title={Before you move on}]
How many English pronouns does \textbf{\gu{તે}} cover? (Three — \emph{he}, \emph{she}, \emph{it}.) What else does the same word do? (Points: \emph{that}.)
\end{tcolorbox}
\section[ame / āpaṇe]{\gu{અમે} / \gu{આપણે} (ame / āpaṇe) — the two we}

\label{lesson:GU-C36-ame-aapne}



\begin{quote}
One lesson back: the pronoun for someone absent. Four chapters back, at the top of its chapter: the word that joins two things.

\begin{itemize}
\raggedright
  \item \emph{Recall:} \emph{te}, and \emph{ane}
\end{itemize}
\end{quote}

\begin{cousinweb}
English has one \emph{we} and it is ambiguous. Gujarati has two, and they are not interchangeable:

\begin{itemize}
\raggedright
  \item \textbf{\gu{અમે}} — \emph{we}, \textbf{not including you}
  \item \textbf{\gu{આપણે}} — \emph{we}, \textbf{including you}
\end{itemize}

If you say \textbf{\gu{અમે}} to someone, you have told them they are outside the group. If you say \textbf{\gu{આપણે}}, you have invited them in. Neither is more polite; they mean different things, and the mistake is audible.

This distinction has a name — linguists call it exclusive and inclusive \emph{we} — and a great many languages make it. English speakers have to learn to hear a choice they have never had to make.
\end{cousinweb}

\subsection*{Your turn}
\begin{itemize}
\raggedright
  \item \emph{Say it:} \emph{ame} — a group the listener is not in.
  \item \emph{Say it:} \emph{āpaṇe} — a group that includes them.
  \item \emph{Contrast:} the two, twice, thinking each time about who is standing there.
\end{itemize}

\begin{tcolorbox}[breakable,colback=teal!4,colframe=teal!35!black,title={Before you move on}]
You are inviting someone to come along. Which \emph{we} do you use, and what would the other one tell them? (\textbf{\gu{આપણે}} — \textbf{\gu{અમે}} would say they are not included.)
\end{tcolorbox}
\section[te]{\gu{તે} reaches the page}

\label{lesson:GU-C36-te-write}



\begin{quote}
Three chapters back you wrote \textbf{\gu{કે}}. Write it once from memory.
\end{quote}

\subsection*{Script — the same stroke, a third time}
\begin{quote}
\textbf{\gu{ત}} — \emph{ta} · \textbf{\gu{ે}} — the \emph{e} stroke
\end{quote}

Set \textbf{\gu{અને}}, \textbf{\gu{કે}} and \textbf{\gu{તે}} in a row and look only at the stroke on top. It is the same \textbf{\gu{ે}} in all three, sitting on three different letters. Learning one stroke has now spelled three of this book's most useful words.

\subsection*{Writing — guided copy, model in view}
Copy \textbf{\gu{તે}} once with the model in view. Then write the row — \textbf{\gu{અને} \gu{કે} \gu{તે}} — and check the three strokes sit at one height.

\begin{tcolorbox}[breakable,colback=teal!4,colframe=teal!35!black,title={Before you move on}]
Write \textbf{\gu{તે}} from memory. Which sign does it share with \textbf{\gu{કે}} and \textbf{\gu{અને}}? (The \textbf{\gu{ે}} above.)
\end{tcolorbox}
\section[je]{\gu{જે} (je) — the one who}

\label{lesson:GU-C36-je}



\begin{quote}
One lesson back you wrote \textbf{\gu{તે}}. One chapter back, the word that places one event against another in time.

\begin{itemize}
\raggedright
  \item \emph{Recall:} \textbf{\gu{તે}}, and \emph{jyāre}
\end{itemize}
\end{quote}

\begin{cousinweb}
\begin{quote}
\textbf{\gu{જે}} — \emph{je} — \textbf{the one who}, \textbf{the thing that}
\end{quote}

Now the stem from \textbf{\gu{જ્યારે}} stands alone. \textbf{\gu{જે}} opens a clause that describes somebody or something, instead of naming them:

\begin{quote}
\textbf{\gu{જે}} \emph{(what they do)} … — \textquotedblleft{}the one who …\textquotedblright{}
\end{quote}

Notice the pairing: \textbf{\gu{જ્યારે}} was \emph{when}, \textbf{\gu{જે}} is \emph{who} or \emph{which}, and both begin the same way. Gujarati marks a describing clause at its front, the way it marked a condition with \textbf{\gu{જો}}.
\end{cousinweb}

\subsection*{Your turn}
\begin{itemize}
\raggedright
  \item \emph{Say it:} \emph{je}, then something a person might do.
  \item \emph{Contrast:} \emph{jyāre} and \emph{je} — the same opening stem, different jobs.
  \item \emph{Build:} describe someone instead of naming them.
\end{itemize}

\begin{tcolorbox}[breakable,colback=teal!4,colframe=teal!35!black,title={Before you move on}]
Describe a person rather than naming them. What does \textbf{\gu{જે}} share with \textbf{\gu{જ્યારે}}? (The opening \textbf{\gu{જ}-} stem; both start their clause.)
\end{tcolorbox}
\section[je … te …]{\gu{જે} … \gu{તે} … (je … te …) — the one who … that one …}

\label{lesson:GU-C36-je-te}



\begin{quote}
The two words this chapter has given you, and the two-part condition from the chapter before.

\begin{itemize}
\raggedright
  \item \emph{Recall:} \emph{je}, \emph{te} — and \emph{jo … to …}
\end{itemize}
\end{quote}

\begin{cousinweb}
\begin{quote}
\textbf{\gu{જે}} \emph{(description)}, \textbf{\gu{તે}} \emph{(the one meant)}
\end{quote}

This is the same shape as \textbf{\gu{જો} … \gu{તો} …}, and recognising that is most of the work. Gujarati marks \textbf{both} ends: the clause that describes, and the thing it describes. English collapses the two — \emph{the one who came is my friend} — and Gujarati keeps them visibly paired.

Once you see the pair you can hear it everywhere: \textbf{\gu{જો} … \gu{તો}}, \textbf{\gu{જ્યારે} … \gu{ત્યારે}}, \textbf{\gu{જે} … \gu{તે}}. One family, one habit, three uses.
\end{cousinweb}

\subsection*{Your turn}
\begin{itemize}
\raggedright
  \item \emph{Build:} \emph{je} + a description, \emph{te} + who is meant.
  \item \emph{Contrast:} the same idea said with a plain name, then with the pair.
  \item \emph{Say it:} \emph{jo … to …} then \emph{je … te …} — hear that they are one shape.
\end{itemize}

\begin{tcolorbox}[breakable,colback=teal!4,colframe=teal!35!black,title={Before you move on}]
Describe someone with the pair. Which earlier pattern is this the same shape as? (\textbf{\gu{જો} … \gu{તો} …}.) Next chapter: the question words, all of them built on one sound.
\end{tcolorbox}
